% !TeX root = ../../Book3.tex
\section{共形度量}
\label{b3:sec:1401}

第十三章的 Poincaré 密度同时具备三种性质：它给出路径长度，在全纯坐标之间按导数的模变换，并且具有常负曲率。本章保留前两种结构，把常曲率条件放宽为一个曲率不等式。为此，先说明一般密度怎样给出距离、拉回在临界点发生什么，以及完备性究竟要求控制哪些路径。

\subsection{共形度量}
\label{b3:subsec:140101}

设 \(U\subseteq\mathbb C\) 是区域。正函数 \(\rho:U\to(0,\infty)\) 定义长度元
\[
 \mathrm ds_\rho=\rho(z)|\mathrm dz|.
\]
它在每一点把所有方向的欧氏长度乘以同一个因子，因而保持夹角，称为共形长度元。其面积元是 \(\rho(z)^2\,\mathrm dA(z)\)，因为两个相互垂直方向的长度都乘以 \(\rho(z)\)。讨论长度与拓扑只需 \(\rho\) 连续；计算曲率时则要求 \(\rho\) 为 \(C^2\) 函数。

\ProseParagraphBreak
若 \(w=F(z)\) 是全纯坐标变换，\(F'\ne0\)，同一长度元在两种坐标中的密度须满足
\[
 \rho_z(z)=\rho_w(F(z))\cdot|F'(z)|.
\]
因此，Riemann 曲面上的一个共形度量，是每张全纯坐标图中的正密度，且这些密度在重叠处满足此式。沿分段 \(C^1\) 路径 \(\gamma\) 定义
\[
 \ell_\rho(\gamma)=\int\rho(\gamma(t))\cdot|\gamma'(t)|\,\mathrm dt.
\]
曲面上使用有限个坐标段分别积分；变换公式和链式法则保证各段的数值不依赖坐标，进一步细分也不改变总长度。以全部连接路径的长度下确界定义 \(\mathrm d_\rho(p,q)\)。对于正密度，这确实给出曲面的距离，但证明正定性时还要考虑路径离开坐标盘的可能。

\Needspace{6\baselineskip}
\begin{proposition}[正密度的路径距离]\label{b3:prop:14-positive-density-distance}
设 \(X\) 是连通 Riemann 曲面，\(\rho\) 是 \(X\) 上正连续共形密度。由分段 \(C^1\) 路径长度的下确界定义的 \(\mathrm d_\rho\) 是有限值距离，其拓扑等于 \(X\) 原有的拓扑。具体地，固定 \(p\in X\)，在一个坐标中令 \(p\) 对应 \(a\)，并取紧包含于该坐标图的闭圆盘 \(\overline{D(a,r)}\)。若该盘上的密度满足 \(0<m\le\rho\le M\)，则对盘内坐标为 \(z\) 的点 \(q\)，有
\[
 m|z-a|\le\mathrm d_\rho(p,q)\le M|z-a|.
\]
对不在开坐标盘 \(D(a,r)\) 内的点 \(q\)，还有 \(\mathrm d_\rho(p,q)\ge mr\)。
\end{proposition}
\begin{proof}
坐标内的直线段给出上界。对于任意连接 \(p,q\) 的路径，若它一直留在该坐标盘中，则密度下界和欧氏曲线长度下界给出 \(\ell_\rho(\gamma)\ge m|z-a|\)。若路径离开该盘，取第一次到达圆周的初始路径段。该段的欧氏长度至少为 \(r\)，故整个路径长度至少为 \(mr\ge m|z-a|\)。这证明盘内点的下界，也同时证明盘外点的下界。

连通曲面上的任意两点可用有限串相交坐标盘及盘内折线路径连接；每段路径的像紧，连续密度在其上有界，所以距离有限。上述局部下界保证不同点的距离为正。常路径、反向路径和路径拼接分别给出零对角线、对称性与三角不等式；三角不等式的下确界步骤与命题~\ref{b3:prop:13-disk-path-metric} 相同，只需分别选取长度在距离的 \(\varepsilon/2\) 以内的两条路径。

欧氏坐标收敛由局部上界推出距离收敛。反之，若 \(\mathrm d_\rho(p,q_n)\to0\)，盘外下界保证 \(q_n\) 最终留在上述坐标盘中，盘内下界再保证其坐标趋于 \(a\)。两种拓扑因此相同。
\end{proof}

拉回全纯映射时还会遇到非负密度。这里称连续非负的共形长度元为共形伪度量，允许它在某些点退化；若其零点离散，则在零点之外仍可按正度量计算。这个用语指长度元允许为零，并不预先保证由它定义的路径下确界能分离各点。后面的曲率比较只使用密度的连续性及其正值处的二阶微分，不在零点定义 \(\log\rho\)。

\subsection{Gaussian 曲率}
\label{b3:subsec:140102}

对正 \(C^2\) 共形密度，本章沿用曲率表达
\[
 K_\rho\coloneqq-\frac{\Delta\log\rho}{\rho^2},
 \quad \Delta=\partial_x^2+\partial_y^2.
\]
它与二维几何中 Gaussian 曲率的认同，按第十三章所说明的范围作为几何输入\footnote{Ahlfors《Conformal Invariants》\cite[1-5 Ultrahyperbolic Metrics, p.~12]{Ahlfors2010Conformal}使用这一表达讨论共形曲率。本书需要的坐标不变性已在命题~\ref{b3:prop:13-curvature-coordinate} 中完整计算；曲率与一般内在几何定义的认同不另假定为本书已证明的曲率张量理论。}。下面恢复本章使用它时的条件与符号。分母记录面积尺度，分子记录密度对数的二阶变化；常密度的曲率因而为零。

\Needspace{6\baselineskip}
\begin{proposition}\label{b3:prop:14-curvature-formula}
设 \(X\) 是 Riemann 曲面，\(\rho\) 是其正 \(C^2\) 共形密度。在任意局部全纯坐标中定义
\[
 K_\rho=-\rho^{-2}\Delta\log\rho,
 \quad \Delta=\partial_x^2+\partial_y^2.
\]
则这一表达给出 \(X\) 上与坐标无关的函数。对常数 \(\kappa>0\)，条件 \(K_\rho\le-\kappa\) 等价于各坐标中均有
\[
 \Delta\log\rho\ge\kappa\rho^2.
\]
\end{proposition}
\begin{proof}
在两张重叠坐标之间，设 \(w=F(z)\)，则 \(F'\ne0\)，密度关系为 \(\rho_z=(\rho_w\circ F)\cdot|F'|\)。命题~\ref{b3:prop:13-curvature-coordinate} 已在正 \(C^2\) 密度及非零全纯坐标导数的条件下证明
\[
 -\frac{\Delta_z\log\rho_z}{\rho_z^2}
 =\left(-\frac{\Delta_w\log\rho_w}{\rho_w^2}\right)\circ F.
\]
该命题使用的两项恒等式是 \(\Delta_z(v\circ F)=|F'|^2(\Delta_wv)\circ F\) 与 \(\Delta\log|F'|=0\)，此处的坐标与密度满足全部条件，故公式直接适用。最后，将 \(-\rho^{-2}\Delta\log\rho\le-\kappa\) 乘以正数 \(\rho^2\) 并移项，得到所述等价不等式。
\end{proof}

对径向函数 \(v=v(r)\)，在 \(r>0\) 处有 \(\Delta v=v''+v'/r\)。将 \(v=\log(1\pm r^2)\) 代入，并在原点按连续性延拓，得到
\[
 \Delta\log\frac{2}{1-|z|^2}
 =\frac{4}{(1-|z|^2)^2},
 \quad
 \Delta\log\frac{2}{1+|z|^2}
 =-\frac{4}{(1+|z|^2)^2}.
\]
连同 \(\Delta\log1=0\)，即得\cref{b3:tab:14-curvature-models} 中的三种曲率。球面密度在平面坐标中趋于零，但这只是该坐标接近无穷远点的表现：换成 \(\zeta=1/z\) 后，密度重新成为 \(2/(1+|\zeta|^2)\)，在 \(\zeta=0\) 处光滑且为正。

\begin{table}[htbp]
\centering\small\renewcommand{\arraystretch}{1.18}\setlength{\tabcolsep}{4pt}
\caption[三种标准共形曲率]{三种标准共形曲率。球面行采用球面去掉一点后的平面坐标；补上无穷远点才得到整张球面。}
\label{b3:tab:14-curvature-models}
\begin{tabularx}{\linewidth}{@{}>{\raggedright\arraybackslash}p{2cm}>{\centering\arraybackslash}p{3.3cm}>{\centering\arraybackslash}p{1.3cm}>{\raggedright\arraybackslash}X@{}}
\toprule
模型 & 密度 \(\rho(z)\) & 曲率 \(K\) & 坐标定义域 \\
\midrule
Euclid 平面 & \(1\) & \(0\) & \(\mathbb C\) \\
Poincaré 圆盘 & \(2/(1-|z|^2)\) & \(-1\) & \(|z|<1\) \\
单位球面 & \(2/(1+|z|^2)\) & \(1\) & \(\mathbb C\)；另一坐标补无穷远点 \\
\bottomrule
\end{tabularx}
\end{table}

若把密度乘以常数 \(c>0\)，则
\[
 \ell_{c\rho}=c\ell_\rho,
 \quad \mathrm d_{c\rho}=c\,\mathrm d_\rho,
 \quad K_{c\rho}=c^{-2}K_\rho.
\]
前两式来自长度积分，第三式来自 \(\log(c\rho)=\log c+\log\rho\)。面积同时乘以 \(c^2\)，而 Cauchy 列与完备性不变。例如，密度 \(1/(1-|z|^2)\) 的曲率是 \(-4\)；将它乘以 \(2\)，才得到本卷使用的曲率 \(-1\) 的圆盘度量。曲率比较中的常数因此必须与长度归一化一起保留。

\ProseParagraphBreak
局部曲率还不能辨认全局拓扑。去心平面上的密度 \(\rho(z)=1/|z|\) 满足 \(\Delta\log\rho=0\)，因为 \(\log|z|\) 在 \(z\ne0\) 处调和。在任一有对数分支的邻域中，\(\mathrm dw=\mathrm dz/z\)，故这个长度元就是 \(w=\log z\) 坐标中的 Euclid 长度元；绕原点一周后 \(w\) 增加 \(2\pi\mathrm i\)，于是全局模型是平直圆柱。平面与圆柱的曲率都为零，闭路和覆盖资料却不同。

\subsection{拉回度量}
\label{b3:subsec:140103}

设 \(f:X\to Y\) 是 Riemann 曲面之间的全纯映射，目标有正 \(C^2\) 共形密度 \(\rho\)。在源坐标 \(z\) 与目标坐标 \(w\) 中，令 \(F=w\circ f\circ z^{-1}\) 是 \(f\) 的局部表达，定义拉回密度
\[
 h_z(z)=\rho_w(F(z))\cdot|F'(z)|.
\]
换目标坐标时，\(\rho_w\) 与 \(F'\) 中的坐标变换导数相消；换源坐标时，\(h_z\) 正好按共形密度的规律变换。因此它定义了 \(X\) 上的全局非负长度元，记为 \(f^*\rho\)。这个记号指拉回后的密度系数，而每个实际计算仍在局部坐标中进行。

\Needspace{6\baselineskip}
\begin{proposition}\label{b3:prop:14-pullback-curvature}
设 \(X,Y\) 是连通 Riemann 曲面，\(f:X\to Y\) 非常值全纯，\(\rho\) 是 \(Y\) 上正 \(C^2\) 共形密度。令 \(h=f^*\rho\)，即局部坐标中的 \(h=(\rho\circ f)\cdot|f'|\)。则 \(h\) 连续，其零点恰为 \(f\) 的临界点，且这些零点离散。在非临界点处，\(h\) 为正 \(C^2\) 密度，并满足
\[
 K_h=K_\rho\circ f.
\]
对任意分段 \(C^1\) 源路径 \(\gamma\)，包括经过临界点的路径，均有
\[
 \ell_h(\gamma)=\ell_\rho(f\circ\gamma).
\]
\end{proposition}
\begin{proof}
局部表达 \(F\) 非常值，故 \(F'\) 不恒为零，其零点离散。目标密度严格为正，所以 \(h\) 的零点正是 \(F'\) 的零点。坐标变换导数不为零，这个零点集合不依赖坐标。连续性由拉回公式得到。

在 \(F'\ne0\) 的点附近，\(F\) 是全纯坐标变换，命题~\ref{b3:prop:14-curvature-formula} 的变换式直接给出 \(K_h=K_\rho\circ f\)。在任意坐标路径段上，链式法则给出
\[
 \rho(F(\gamma(t)))\cdot|(F\circ\gamma)'(t)|
 =\rho(F(\gamma(t)))\cdot|F'(\gamma(t))|\cdot|\gamma'(t)|.
\]
这个等式在 \(F'=0\) 时仍成立；积分并将有限个坐标段相加，便得长度等式。
\end{proof}

临界点附近的退化有明确形式。若 \(F'\) 在 \(a\) 处有 \(m\) 阶零点，写成 \(F'(z)=(z-a)^m g(z)\)，其中 \(g\) 全纯且在邻域内无零点，则
\[
 h(z)=|z-a|^m b(z),
 \quad b(z)=\rho(F(z))\cdot|g(z)|>0,
\]
且 \(b\) 为 \(C^2\) 函数。因子 \(|z-a|^m\) 说明拉回密度在临界点可能不再是正 \(C^2\) 密度；曲率公式不能在那里除以 \(h^2\)。若 \(f\) 为常值，整个拉回密度都为零。下一节将允许这些退化，并只在密度为正的比较点计算曲率。

\subsection{完备性}
\label{b3:subsec:140104}

正连续密度 \(\rho\) 称为完备的，是指其路径距离 \(\mathrm d_\rho\) 的每个 Cauchy 列都在曲面内收敛。正性保证局部距离与原拓扑一致，完备性则限制路径能否用有限长度离开曲面的全部紧部分。为了精确表达后一点，称连续路径 \(\gamma:[0,1)\to X\) 逃离紧集，是指对任意紧集 \(K\subseteq X\)，都存在 \(t_K<1\)，使 \(\gamma(t)\notin K\) 对所有 \(t>t_K\) 成立。

\Needspace{6\baselineskip}
\begin{proposition}[完备性的逃逸路径判据]\label{b3:prop:14-completeness-escape}
设 \(X\) 是连通 Riemann 曲面，\(\rho\) 是其正连续共形密度，\(\mathrm d_\rho\) 是相应路径距离。则 \((X,\mathrm d_\rho)\) 完备，当且仅当每条逃离紧集的连续、局部分段 \(C^1\) 路径 \(\gamma:[0,1)\to X\) 都具有无限长度。这里总长度定义为
\[
 \ell_\rho(\gamma)=\sup_{0<t<1}\ell_\rho(\gamma|_{[0,t]}).
\]
\end{proposition}
\begin{proof}
先设距离完备。若有逃离紧集的路径长度有限，则其尾部长度趋于零。取 \(t_n\uparrow1\)，对 \(m>n\) 有
\[
 \mathrm d_\rho(\gamma(t_n),\gamma(t_m))
 \le\ell_\rho(\gamma|_{[t_n,t_m]})\longrightarrow0
\]
且上界对所有 \(m>n\) 一致。因此 \(\gamma(t_n)\) 是 Cauchy 列，收敛到某个 \(p\in X\)。由距离拓扑与曲面拓扑相同，取 \(p\) 的一个紧坐标邻域后，该列最终留在此紧集中，与路径逃离每个紧集矛盾。

反之，若距离不完备，取一个没有极限的 Cauchy 列。可抽取子列 \((p_n)\)，使
\[
 \mathrm d_\rho(p_n,p_{n+1})<2^{-n-2}.
\]
按路径下确界的定义，选取连接 \(p_n,p_{n+1}\) 的分段 \(C^1\) 路径 \(\gamma_n\)，使 \(\ell_\rho(\gamma_n)<2^{-n-1}\)。把这些路径依次线性改参数到区间 \([1-2^{1-n},1-2^{-n}]\)，拼成 \(\gamma:[0,1)\to X\)。它连续且局部分段 \(C^1\)，总长度不超过 \(\sum_{n\ge1}2^{-n-1}<\infty\)。

还须证明这条路径确实逃离紧集。若某个紧集 \(K\) 被任意靠后的路径段访问，就能选取 \(n_j\to\infty\) 和 \(q_j\in K\cap\gamma_{n_j}\)。由紧性及距离拓扑与原拓扑相同，可再取子列使 \(q_j\to q\in K\)。而
\[
 \mathrm d_\rho(p_{n_j},q_j)
 \le\ell_\rho(\gamma_{n_j})<2^{-n_j-1}\longrightarrow0,
\]
故 \(p_{n_j}\to q\)。原来的 Cauchy 列一旦有收敛子列，就必须整体收敛到同一点，矛盾。因此 \(\gamma\) 逃离每个紧集，构成有限长度的逃逸路径。这证明逆命题。
\end{proof}

这个判据也给出实用的比较办法。若 \(\sigma\) 是同一曲面上的完备正连续密度，而 \(\rho\ge c\sigma\)、\(c>0\)，则每条逃逸路径都有
\(\ell_\rho\ge c\ell_\sigma=\infty\)，所以 \(\rho\) 完备。例如，平面密度 \(e^{|z|^2}\ge1\)，由 Euclid 平面的完备性即可知道它完备；无需先寻找该度量的最短路径。

\ProseParagraphBreak
单位盘的常密度 \(1\) 不完备，因为径向路径可在有限长度内趋近单位圆周。圆盘 Poincaré 密度的完备性已在第十三章证明；若将同一密度限制到 \(D(0,a)\)、\(0<a<1\)，曲率仍恒为 \(-1\)，但趋近实轴端点 \(a\) 的径向长度只有 \(2\operatorname{arctanh}a\)，所以这个受限度量不完备。负曲率条件由局部微分表达控制，端部是否在有限距离处还须另行判断。

\ProseParagraphBreak
表中的球面密度也说明不能忘记曲面本身的范围。在 \(\mathbb C\) 上，\(\rho(z)=2/(1+|z|^2)\) 沿正实轴到无穷远的总长度为 \(\int_0^\infty2/(1+t^2)\,\mathrm dt=\pi\)，故去掉无穷远点的度量不完备；补上该点后，密度在另一坐标中仍正且光滑，所得球面是紧度量空间，因而完备。

\ProseParagraphBreak
应用逃逸判据时，一条有限长度的逃逸路径就足以否定完备性；证明完备则必须控制所有逃逸路径。只算几条径向线的长度不能承担后一种量词。下一节的 Ahlfors--Schwarz 比较直接使用局部曲率和小圆盘最大值，不把完备性列为待比较密度的假设。
